\section{Interpretation of Research Questions}
\label{sec:discussion-rqs}

\begin{table}[ht]
\centering
\caption{Current evidential status of the research questions.}
\label{tab:rq-status}
\begin{tabular}{|C{1.3cm}|L{4.4cm}|L{5.3cm}|L{2.8cm}|}
\hline
\textbf{RQ} & \textbf{Objective} & \textbf{Available evidence} & \textbf{Status} \\
\hline
RQ1 & Decision conformance, enforcement correctness, evidence completeness &
$106/106$ pre-registered runs conformant; six enforcement invariants held on
every applicable run; ten degradation scenarios behaved as predicted &
Evaluated \\
\hline
RQ2 & Equivalent governance of an on-premises target and a cloud target by the
shared risk engine & All seven on-premises checkpoints attained, the last three
in a $6$-run live-target arm; of eleven catalogued weakness classes, three named
on both sides of the boundary at the frozen gate commit and ten after
remediation, in a shared finding vocabulary &
Evaluated (coverage bounded) \\
\hline
RQ3 & Multi-tool risk estimation and deployment decision & Held-out model results
available; complete gate-decision analysis incomplete & Partially evaluated \\
\hline
\end{tabular}
\end{table}

\subsection{RQ1: Decision Conformance and Enforcement Correctness}
\label{subsec:discussion-rq1}

RQ1 asks three things: whether the decision reached is the decision the case
warranted, whether that decision constrains what happens next, and whether the
evidence behind it is preserved. The campaign answers all three positively
(Tables~\ref{tab:rq1-confusion} and~\ref{tab:rq1-invariants}): the confusion
matrix is exactly diagonal, no \emph{reject} was followed by a deployment, and
no \emph{review} was deployed without a matching approval record.

What this establishes is internal correctness under pre-declared conditions. Two things
limit it. The fixtures were authored by the same work that built the gate, so
the result shows that the implementation satisfies its own specification rather
than that the specification assigns the right band to arbitrary real
infrastructure. And the invariants were not exercised equally often: the
approval-record invariant rests on ten runs and the rejection invariants on
sixteen. Ten clean runs show that the mechanism works when it is invoked; they
do not show how often it would fail.

\subsubsection{The Gate Fails Open}

The most consequential result in the RQ1 data is not the conformance figure but
the degradation table. Removing Checkov from a \emph{review}-band cloud change
drops its score from $54$ to $9$ and releases it as \emph{pass}
(Table~\ref{tab:rq1-degradation}). The gate conformed to its specification here,
since the scenario declared that outcome in advance. The problem is that the
specification permits the release: the evidence that would have held the change
was never collected.

This follows from the additive form of the scoring function. In
$S = \sum_{f} w(\mathrm{sev}(f)) + c(\Delta \mathrm{cost}) + 5v + m$, a scanner
that was never invoked and a scanner that ran and found nothing both contribute
zero. The score therefore conflates \emph{absence of evidence} with
\emph{evidence of absence}, the classic failure mode of any evidence-summing
control. The system does record which scanners were unavailable, so an auditor
can detect the condition after the fact, but nothing in the decision function
consumes that field. The record is diagnostic rather than protective.

Chapter~\ref{Chapter6}
proposes an assurance penalty --- a positive contribution for each unavailable
scanner, or a fail-closed policy above a coverage threshold --- as the remedy.

\subsubsection{Gate Overhead}

Both branches are acceptable in absolute terms for a pre-deployment check, so
the concern is the scaling behaviour rather than the current figure. The cloud
gate consumes roughly seven-tenths of run wall-clock time, and its cost is
dominated by Checkov being invoked once per synthesised template file. It
therefore grows with the size of the deployed stack rather than with the size of
the change under review. A production deployment with many templates would find
the gate, not the cloud provider, to be the bottleneck.

\subsection{RQ2: Cross-Boundary Enforcement Equivalence}
\label{subsec:discussion-rq2}

RQ2 asks whether the two sides of the boundary are governed \emph{alike}.

The workflow half is answered positively across all seven checkpoints. The last
three --- apply, post-apply verification, and idempotency --- were closed by a live-target arm. That arm supplies a kind of evidence
the offline campaign structurally could not: an applied state read back and
asserted against independently written expectations, a second apply that reports
\texttt{changed=0}, and three runs that requested deployment of the rejected
fixture against a reachable, authenticated host and were stopped at the gate.
Enforcement holds at the point of application, not merely in simulation.

The detection half produced the more interesting finding, and it is not a
finding about decisions. At the frozen commit only three of the eleven
catalogued classes were named on both branches
(Table~\ref{tab:rq2-coverage}). A gate that treats the branches as
interchangeable would have been wrong about eight of eleven classes, and nothing
in the architecture would have revealed it. Where a class \emph{was} named on
both sides, the identifiers matched: an unrestricted firewall rule is reported as
\texttt{sg\_ssh\_open} whether it appears in a security group or in an
\texttt{iptables} task, and unrestricted administrative privilege as
\texttt{iam\_wildcard} whether it comes from an IAM policy or a sudoers file.
That shared vocabulary is the substantive achievement, because it is what makes
the two sides comparable at all. Equal scores would not have been.

A shared scoring
function and shared thresholds are frequently taken to imply consistent
treatment across a hybrid estate. They do not. Consistency of \emph{decision
logic} is worth little without comparable \emph{detection coverage} feeding it,
and the boundary at which coverage diverges is invisible in any architecture
diagram --- both branches appear to call the same gate, because they do. Nor is
such a gap reachable by tuning: a weakness that produces no finding cannot be
scored, reviewed, or audited at any threshold, so the remedy is to add the
missing detection rather than to reweight what is already present. The
remediated arm shows that this remedy was available and inexpensive, at nine
rule families in two files. It should not be read as a recall estimate: it
measures whether the gaps the baseline located were closable within this
architecture, not how much the gate sees in general.

The bound on this evidence is no longer the size of the class list but the
authorship of the fixtures. Eleven classes drawn from an external catalogue
remove the author's discretion over which weaknesses are counted, but not over
how each one is written, so the figures measure rule-to-CWE alignment rather
than field recall. The specificity controls narrow this from the other side, but
a control that produces no findings at all cannot establish that a matcher will
refrain from firing on an unrelated finding in a busier report --- a failure
mode an earlier candidate control exhibited before it was replaced.

\subsection{RQ3: Unified Risk-Evaluation Capability}
\label{subsec:discussion-rq3}

The classifier experiment provides positive but bounded evidence for RQ3. The
logistic-regression component achieved an accuracy of $0.9259$, an F1-score of $0.9240$,
and a ROC--AUC of $0.9233$ on the held-out partition. The confusion matrix contained 8
false positives and 16 false negatives. Within the filtered corpus, the combined Bandit and Semgrep features therefore
provided a useful learned risk signal.

This result does not validate the complete gate. The model was evaluated only on
detector-visible Python vulnerabilities, and categories outside the capabilities of the
feature extractors were excluded. In addition, the negative files were sampled from
maintained projects rather than formally verified as secure. The controlled tests for
score aggregation, threshold boundaries, deduplication, and final deployment decisions
also remain incomplete. RQ3 is consequently interpreted as partially evaluated: the
learned component performed consistently on the selected test set, but the complete
multi-component decision mechanism has not yet been fully validated.

\section{Comparison with Related Approaches}
\label{sec:discussion-comparison}

The proposed process combines cloud operations and continuous security in one governed
workflow. Table~\ref{tab:discussion-comparison} summarises its relationship to established
approaches without claiming quantitative superiority, since no manual or ungated baseline
was evaluated.

\begin{table}[ht]
\centering
\small
\setlength{\tabcolsep}{4pt}
\caption{Concise comparison with related cloud operational approaches.}
\label{tab:discussion-comparison}
\begin{tabular}{|L{3.2cm}|L{4.4cm}|L{6.3cm}|}
\hline
\textbf{Approach} & \textbf{Primary emphasis} & \textbf{Relationship to this research} \\
\hline
Traditional operations & Separately coordinated provisioning, administration, and
review & Integrates these stages through one gate and audit trail; operational gains
remain unmeasured. \\
\hline
NIST reference architecture \cite{Liu2011NISTRA} & Vendor-neutral actors and cloud
components & Adds an executable operational sequence, but is less vendor-neutral. \\
\hline
CloudCAMP and IaC practice \cite{Bhattacharjee2019CloudCAMP,Guerriero2019IaC,Hasan2020IaCTesting}
& Automated provisioning, abstraction, and IaC quality & Adds gating and traceability,
while retaining adapter-maintenance costs. \\
\hline
Continuous security \cite{Kumar2020ADOC,Rajapakse2021ToolIntegration} & Lifecycle security,
tool integration, and triage & Adds hybrid-cloud decision control and a shared finding
vocabulary; source evidence remains necessary. \\
\hline
\end{tabular}
\end{table}

The contribution is therefore an operational integration of existing ideas: it is more
concrete than a reference architecture and more governance-oriented than provisioning
automation, but less portable and not yet proven more efficient than traditional practice.

\section{Limitations}
\label{sec:discussion-limitations}

The study has seven principal limitations.

\begin{enumerate}
    \item \textbf{Self-authored fixtures and limited scale.} RQ1 and RQ2 were
    evaluated on a pre-registered campaign of 106 runs whose scenarios and
    expected outcomes were fixed in advance, which removes the selection
    freedom of an ad-hoc development sample. It does not remove a more basic
    dependency: the fixtures were authored by the same work that built the
    gate. Conformance therefore establishes that the implementation satisfies
    its own specification on cases whose correct treatment is already known,
    not that the specification assigns the right band to arbitrary real
    infrastructure. Establishing the latter would require independently
    authored cases. The prototype has
    additionally not been evaluated with large repositories, concurrent
    pipelines, or a substantial private-infrastructure fleet.

    \item \textbf{Narrow deployment-side evidence.} Apply, post-apply
    verification, and repeated-run idempotency were measured only in the
    live-target arm: three replicates of one scenario against a single
    disposable virtual machine running one operating system image. That
    supports two claims --- that the deployment path completes correctly, and
    that a \emph{reject} decision blocks application against a real host. It
    does not support claims about host heterogeneity, concurrent targets,
    partial failures, or hosts whose prior state differs from a fresh image,
    which are exactly the conditions under which convergence tooling tends to
    misbehave.

    \item \textbf{Restricted implementation scope.} The learned component is limited to
    Python and uses features from Bandit and Semgrep. The wider prototype relies on a
    small fixed set of adapters and is centred on AWS CDK, CloudFormation, AWS operational
    services, and Ansible. The findings cannot be assumed to transfer to other languages,
    cloud providers, configuration systems, or vulnerability classes.

    \item \textbf{Risk assumptions and the fail-open score.} The severity weights,
    cost bands, compliance contribution, model contribution, and decision
    thresholds have not been calibrated using expert judgement or historical
    incidents. Exposure, asset criticality, data sensitivity, and compensating
    controls are represented only indirectly. More seriously, the campaign
    confirmed that the additive score \emph{fails open}: an unavailable scanner
    and a scanner that found nothing contribute identically
    (Table~\ref{tab:rq1-degradation}). The condition is recorded in the
    scanner-status field, but no part of the decision function acts on it.

    \item \textbf{Unequal detection coverage across the boundary.} The two
    branches share a decision function but not a scanner set. Coverage is also established over hand-authored
    fixtures, so the figures bound rule-to-CWE alignment rather than recall on
    real infrastructure, and the specificity controls cannot exclude a matcher
    firing on an unrelated finding in a busier report.

    \item \textbf{Limited comparative and operational measures.} The study does not yet
    compare the proposed process with manual cloud operations or IaC deployment without
    the unified gate. It also does not measure operator effort, deployment lead time,
    change-failure rate, remediation time, false-alert handling, or long-term incident
    reduction. The current findings therefore establish limited component feasibility
    rather than overall operational superiority.
\end{enumerate}

\section{Directions for Improvement}
\label{sec:discussion-improvements}

The improvements follow directly from the observed limitations.

\begin{enumerate}
    \item \textbf{Fail-open scoring:} make assessment availability part of the decision
    by applying an assurance penalty or forcing review below a coverage threshold.

    \item \textbf{Unequal detection coverage:} add source-appropriate detection for the
    remaining cloud-side secret gap and explicitly identify components that do not apply
    to one branch.

    \item \textbf{Limited scale and external validity:} repeat the campaign with
    independently authored fixtures, heterogeneous hosts, concurrent targets, induced
    failures, larger repositories, and incremental scanning.

    \item \textbf{Restricted learned model:} use larger multi-project datasets,
    project-aware partitions, repeated cross-validation, external test sets, and
    language-specific models.

    \item \textbf{Uncalibrated operational value:} calibrate score weights with expert
    and incident data, then compare the process with manual operations and ungated IaC
    using lead time, effort, policy violations, and remediation time.
\end{enumerate}